%% ma: L11-B21-0019
%% lop: 11
%% chuong: 6
%% bai: 21
%% dang: 11.21.1
%% loai: TN4
%% muc: TH
%% dap_an: "C"
%% nguon: "(NBV)-ĐỀ SỐ 19-MỖI NGÀY 1 ĐỀ THI 2025.pdf (D19, MỖI NGÀY 1 ĐỀ - Đề số 19), Phần I Câu 7"
%% kien_thuc: "Bài 21 (bất phương trình lôgarit cơ bản, cơ số lớn hơn 1), bất phương trình bậc hai"
%% trang_thai: chinh-thuc
%% ghi_chu: "Đáp án gốc C đúng (S = (-1;2), b - a = 3). Trùng ý tưởng với dạng đã duyệt 11.21.1 (bất phương trình lôgarit cơ bản), khác ở chỗ biểu thức trong lôgarit là tam thức bậc hai; nhập chính thức."
\begin{ex}
Bất phương trình $\log_3\left(x^2-x+7\right)<2$ có tập nghiệm là $S=(a;b)$. Tính $b-a$.
\choice{$b-a=-1$}{$b-a=-3$}{\True $b-a=3$}{$b-a=1$}
\loigiai{Vì $x^2-x+7=\left(x-\dfrac12\right)^2+\dfrac{27}4>0$ với mọi $x$ nên điều kiện luôn thỏa. Cơ số $3>1$ nên bất phương trình tương đương $x^2-x+7<3^2=9$, tức $x^2-x-2<0\Leftrightarrow-1<x<2$. Vậy $S=(-1;2)$, $b-a=2-(-1)=3$. Chọn C.}
\end{ex}
